% A six-cycle: every vertex has a neighbor in each of two coordinates.
\begin{tikzpicture}[figure picture,foundation picture,x=1cm,y=1cm,
 vertex/.style={draw=FigureGrid,fill=FigurePaper,minimum width=16mm,minimum height=7mm,font=\figurefont\footnotesize},
 first/.style={FigureBlue,line width=1.1pt},second/.style={FigureGreen,line width=1.1pt,dashed}]
 \node[vertex] (v1) at (0,1.8) {$(a,b)$};
 \node[vertex] (v2) at (1.9,0.9) {$(b,b)$};
 \node[vertex] (v3) at (1.9,-0.9) {$(b,c)$};
 \node[vertex] (v4) at (0,-1.8) {$(c,c)$};
 \node[vertex] (v5) at (-1.9,-0.9) {$(c,a)$};
 \node[vertex] (v6) at (-1.9,0.9) {$(a,a)$};
 \draw[first] (v1)--(v2); \draw[second] (v2)--(v3);
 \draw[first] (v3)--(v4); \draw[second] (v4)--(v5);
 \draw[first] (v5)--(v6); \draw[second] (v6)--(v1);
 \draw[first] (3.2,1.5)--(4,1.5);
 \node[anchor=west,font=\figurefont\footnotesize] at (4.1,1.5) {只改$x_1$的预测};
 \draw[second] (3.2,0.8)--(4,0.8);
 \node[anchor=west,font=\figurefont\footnotesize] at (4.1,0.8) {只改$x_2$的预测};
 \node[align=center,font=\figurefont\footnotesize] at (4.6,-0.45)
 {$d_{DS}=2$\\$d_N=1$};
\end{tikzpicture}%
